\chapter{Buổi 1}
\label{Chapter1}
\setcounter{secnumdepth}{5}

\section*{Câu 1}
\addcontentsline{toc}{section}{Câu 1}
Nêu một số ứng dụng của xử lý ảnh số và video số trong việc thông minh hoá quy trình sản xuất trong công nghiệp/thiết kế sản phẩm/thu hoạch sản phẩm/sản xuất/phân phối sản phẩm.

\begin{itemize}
    \item \textbf{Thiết kế sản phẩm:} Quét 3D hoặc camera độ sâu giúp tạo mô hình CAD từ mẫu thật; video tốc độ cao hỗ trợ phân tích biến dạng, va đập hoặc luồng khí quanh sản phẩm.
    \item \textbf{Thông minh hoá quy trình sản xuất:} Camera AI kiểm tra lỗi bề mặt, thiếu linh kiện, sai màu, sai kích thước; đồng thời giám sát an toàn lao động và hỗ trợ robot lắp đặt chính xác.
    \item \textbf{Thu hoạch sản phẩm:} Camera RGB hoặc đa phổ giúp robot nhận biết độ chín, phát hiện sâu bệnh và phân loại nông sản theo kích thước, màu sắc, hình dạng.
    \item \textbf{Phân phối sản phẩm:} OCR, barcode/QR và nhận diện biển số hỗ trợ phân loại hàng, kiểm kho, theo dõi băng chuyền và điều phối xe ra vào kho.
\end{itemize}

\section*{Câu 2}
\addcontentsline{toc}{section}{Câu 2}
Nêu một số ứng dụng của xử lý ảnh số và video số trong quản lý xã hội: eKYC, thương mại điện tử, giám sát an ninh.

\begin{itemize}
    \item \textbf{eKYC:} OCR trích xuất thông tin từ giấy tờ, nhận diện khuôn mặt đối chiếu với ảnh giấy tờ và kiểm tra sống thật bằng video hoặc chuyển động ngẫu nhiên.
    \item \textbf{Thương mại điện tử:} Tìm kiếm bằng hình ảnh, thử sản phẩm bằng AR, gắn nhãn tự động và phát hiện nội dung vi phạm.
    \item \textbf{Giám sát an ninh:} Camera AI nhận diện khuôn mặt, biển số, hành vi bất thường, vi phạm giao thông hoặc dấu hiệu cháy nổ để cảnh báo sớm.
\end{itemize}

\section*{Câu 3}
\addcontentsline{toc}{section}{Câu 3}
Liệt kê các lĩnh vực mà thị giác máy tính mang lại doanh thu lớn.

Các lĩnh vực tạo doanh thu lớn cho thị giác máy tính gồm:
\begin{itemize}
    \item \textbf{Ô tô và xe tự hành:} ADAS, nhận diện làn đường, vật cản, biển báo và giám sát tài xế.
    \item \textbf{Y tế:} Phân tích X-quang, CT, MRI để hỗ trợ phát hiện tổn thương, khối u hoặc bất thường bệnh lý.
    \item \textbf{Sản xuất:} Kiểm tra chất lượng tự động, phát hiện lỗi linh kiện và hướng dẫn robot trong nhà máy thông minh.
    \item \textbf{Bán lẻ và thương mại điện tử:} Cửa hàng không thu ngân, kiểm kho tự động, visual search và thử đồ ảo.
    \item \textbf{An ninh và thành phố thông minh:} Nhận diện khuôn mặt, ANPR, kiểm soát ra vào và phân tích hành vi đám đông.
\end{itemize}

\section*{Câu 4}
\addcontentsline{toc}{section}{Câu 4}
Cho ví dụ về hệ thống thị giác xuất ra thông tin ở mức wisdom.

Ví dụ: hệ thống thị giác của xe tự hành trong tình huống khẩn cấp. Camera phát hiện phía trước có trẻ em, bên cạnh là hàng rào mềm, trong xe có hành khách dễ tổn thương. Thay vì chỉ chọn quỹ đạo theo công thức tối ưu khoảng cách, hệ thống đánh giá bối cảnh, mức độ nguy hiểm và giá trị an toàn con người.

Quyết định ở mức \textbf{wisdom} là chọn hành động giảm thiểu thiệt hại tổng thể: phanh khẩn cấp, đánh lái về phía vật cản ít nguy hiểm hơn, kích hoạt túi khí và siết dây an toàn. Như vậy, hệ thống không chỉ nhận biết vật thể mà còn ra quyết định có xét đến đạo đức, rủi ro và bối cảnh.

\section*{Câu 5}
\addcontentsline{toc}{section}{Câu 5}
Tự đề ra mục đích khi nghiên cứu thị giác máy tính.

Mục đích nghiên cứu của em là hiểu cách máy tính thu nhận, xử lý và diễn giải hình ảnh/video để xây dựng các hệ thống hỗ trợ con người trong sản xuất, y tế, giao thông và đời sống. Trọng tâm là phát triển mô hình có độ chính xác cao, dễ triển khai và có khả năng hoạt động ổn định trong môi trường thực tế.

\section*{Câu 6}
\addcontentsline{toc}{section}{Câu 6}
Tự đề ra mục tiêu cần thực hiện để đặt được mục đích trên.

\begin{itemize}
    \item Nắm vững kiến thức nền tảng về ảnh số, lọc ảnh, trích xuất đặc trưng và học máy.
    \item Thực hành các bài toán cơ bản như phân loại ảnh, phát hiện vật thể, phân đoạn ảnh và theo dõi đối tượng.
    \item Học cách xây dựng, huấn luyện, đánh giá và tối ưu mô hình deep learning cho dữ liệu hình ảnh.
    \item Tìm hiểu các vấn đề triển khai thực tế như dữ liệu nhiễu, thay đổi ánh sáng, tốc độ xử lý và độ tin cậy.
\end{itemize}

\section*{Câu 7}
\addcontentsline{toc}{section}{Câu 7}
Cho ví dụ về một hệ thống thị giác có khả năng thích ứng với sự thay đổi.

Ví dụ: camera giám sát ngoài trời dùng AI và học liên tục. Hệ thống phải hoạt động qua ngày/đêm, mưa, sương mù, thay đổi ánh sáng hoặc thay đổi bối cảnh như đường đang sửa chữa.

Khả năng thích ứng thể hiện ở việc hệ thống tự cân bằng sáng, khử nhiễu, chuyển sang ảnh hồng ngoại khi thiếu sáng, cập nhật mô hình để giảm báo động giả và học thêm các bối cảnh mới. Nhờ vậy, hệ thống vẫn theo dõi được người, xe hoặc hành vi bất thường dù môi trường thay đổi.

\section*{Câu 8}
\addcontentsline{toc}{section}{Câu 8}
Cho ví dụ về một hệ thống thị giác không có khả năng thích ứng với sự thay đổi.

Ví dụ: hệ thống kiểm tra lỗi bo mạch bằng \textbf{template matching}. Hệ thống so sánh ảnh hiện tại với ảnh mẫu cố định nên phụ thuộc mạnh vào ánh sáng, góc chụp và vị trí bo mạch.

Nếu đèn yếu đi, ánh sáng ngoài lọt vào, bo mạch lệch vài độ hoặc khoảng cách camera thay đổi, ảnh không còn khớp với mẫu ban đầu. Hệ thống có thể báo lỗi giả hoặc dừng dây chuyền dù sản phẩm không hỏng. Đây là ví dụ của hệ thống thiếu khả năng thích ứng.

\section*{Câu 9}
\addcontentsline{toc}{section}{Câu 9}
Phân biệt lập trình truyền thống và học máy.

\small
\begin{longtable}{|p{0.22\textwidth}|p{0.34\textwidth}|p{0.34\textwidth}|}
\hline
\textbf{Tiêu chí} & \textbf{Lập trình truyền thống} & \textbf{Học máy} \\
\hline
\textbf{Cách tiếp cận} & Con người viết luật xử lý rõ ràng. & Mô hình học quy luật từ dữ liệu. \\
\hline
\textbf{Đầu vào chính} & Quy tắc và dữ liệu đầu vào. & Dữ liệu huấn luyện và nhãn/feedback. \\
\hline
\textbf{Vai trò con người} & Thiết kế từng bước xử lý. & Chuẩn bị dữ liệu, chọn mô hình, huấn luyện và đánh giá. \\
\hline
\textbf{Khả năng mở rộng} & Khó khi có nhiều ngoại lệ. & Tốt hơn nếu có dữ liệu đa dạng và chất lượng. \\
\hline
\textbf{Ví dụ} & Tính hóa đơn theo công thức cố định. & Nhận diện ảnh, dự đoán giá, phát hiện thư rác. \\
\hline
\end{longtable}
\normalsize

\section*{Câu 10}
\addcontentsline{toc}{section}{Câu 10}
Cho ví dụ về meta learning.

Meta-learning là hướng học máy trong đó mô hình học cách thích ứng nhanh với tác vụ mới dựa trên kinh nghiệm từ nhiều tác vụ trước đó. Nó đặc biệt hữu ích trong bài toán ít dữ liệu.

Ý tưởng này tương tự cách con người học: nếu một người đã biết đi xe đạp, họ thường học đi xe máy nhanh hơn người chưa từng quen với phương tiện hai bánh. Kinh nghiệm trước đó giúp quá trình học nhiệm vụ mới diễn ra nhanh hơn.

Ví dụ: hệ thống nhận diện động vật quý hiếm chỉ có vài ảnh mẫu cho mỗi loài. Trong giai đoạn meta-training, mô hình được học qua nhiều tác vụ phân loại nhỏ như phân biệt chó/mèo, chim/sóc, hổ/báo. Nhờ đó, mô hình học được cách tìm đặc trưng phân biệt tổng quát. Khi gặp một loài mới chỉ có vài ảnh, mô hình vẫn có thể thích ứng nhanh và phân loại tốt hơn so với huấn luyện từ đầu.

Một ví dụ khác là robot nhà máy được huấn luyện trên nhiều nhiệm vụ như gắp, đặt, vặn, kéo. Khi gặp nhiệm vụ mới, robot chỉ cần ít lần thử để điều chỉnh lực và quỹ đạo chuyển động phù hợp.
